\documentclass[11pt]{article}
\usepackage[margin=25mm]{geometry}
\usepackage{amsmath,amssymb,amsthm}
\newtheorem{proposition}{Proposition}
\title{Pair counts across consecutive distance layers}
\author{Human Compression project}
\date{}
\begin{document}\maketitle

Let $G$ be a connected simple four-regular graph of girth at least five,
and let $V_0,\ldots,V_T$ be its distance layers from a specified vertex.
Write $n_i=|V_i|$, $e_i=e(G[V_i])$ and $c_i=e(V_i,V_{i+1})$, with
$c_{-1}=c_T=0$. Thus
\[
 4n_i=2e_i+c_{i-1}+c_i,\qquad n_0=1,\quad n_1=4,\quad n_2=12.
\]
The last three equalities require $T\geq2$. They follow from the absence
of triangles and quadrangles in the radius-two ball. Every nonempty
proper distance prefix has a positive even edge boundary. Consequently
$c_i$ is positive and even for $0\leq i<T$. Each vertex outside $V_0$
has a parent, so $c_{i-1}\geq n_i$ and $c_i\leq3n_i$ for $i>0$.

For $s\geq1$ and $0\leq x\leq4s$, put
\[
 \Phi_s(x)=(s-r)\binom q2+r\binom{q+1}2,
 \qquad x=qs+r,\quad 0\leq r<s.
\]
This is the minimum of $\sum_{v=1}^s\binom{x_v}2$ over nonnegative
integer entries of sum $x$. Moving one unit from an entry at least two
larger than another strictly decreases the sum. A minimum therefore has
only the two values $q,q+1$. Equivalently,
\[
 \Phi_s(x)=\max_{0\leq a\leq3}\left(ax-\binom{a+1}2s\right).
 \tag{1}
\]
For $x<4s$ the maximum is attained at $a=\lfloor x/s\rfloor$;
at $x=4s$ the value at $a=3$ is $6s$, also the balanced minimum.

\begin{proposition}
For $i<j$, define
\[
 E_{ij}=\sum_{h=i}^j e_h+\sum_{h=i}^{j-1}c_h.
\]
Then
\begin{align*}
 E_{ij}&+\Phi_{n_i}(2e_i+c_i)
       +\Phi_{n_j}(2e_j+c_{j-1})+6\sum_{h=i+1}^{j-1}n_h\\
 &\quad+\mathbf1_{i>0}\Phi_{n_{i-1}}(c_{i-1})
       +\mathbf1_{j<T}\Phi_{n_{j+1}}(c_j)
 \leq \binom{\sum_{h=i}^j n_h}{2}.                         \tag{2}
\end{align*}
Terms involving nonexistent layers are omitted. For $i=j$, the two
endpoint terms and the interior sum are replaced by $\Phi_{n_i}(2e_i)$.
\end{proposition}
\begin{proof}
Apply the elementary common-neighbor pair inequality
\[
 e(G[S])+\sum_{v\in V(G)}\binom{|N(v)\cap S|}{2}\leq\binom{|S|}{2}
\]
to $S=V_i\cup\cdots\cup V_j$. An adjacent pair in $S$ has no common
neighbor, since triangles are forbidden; any other pair has at most one,
since quadrangles are forbidden. This proves the displayed inequality.
Its first term is $E_{ij}$. Every vertex in a strictly interior layer has
all four neighbors in $S$, contributing six. The sums of the degrees
into $S$ on the two endpoint layers are $2e_i+c_i$ and
$2e_j+c_{j-1}$. The sums on the two external neighboring layers are
$c_{i-1}$ and $c_j$. No other vertex has a neighbor in $S$.
Balancing within each of these four layers gives (2). When $i=j$,
there is a single internal contribution with degree sum $2e_i$.
\end{proof}

\paragraph{Application and its limits.}
For the existing F009 target, $T=22$ and $\sum n_i\leq94$.
The acquired interval bounds are $(1,2,5,8,17)$ for intervals of one
through five layers; the endpoint bounds are $(1,5,17)$.
Their partition recurrence gives the lower bound $Q_r$ for either end
of $r$ layers, hence $n_i\leq94-Q_i-Q_{22-i}$. At the initial end,
the endpoint argument uses the radius-two ball about the root; at the
terminal end it uses the ball about a vertex of distance exactly $22$.
No assumption that every pair of vertices is at distance at most $22$
is needed. These are the attributed Moore bounds used in
\texttt{corpus106\_F009\_suffix\_reserve\_v1/suffix\_bound.tex}.

The adjacent program combines those bounds, the degree identities,
even cuts, elementary simple-graph capacities and (2) for intervals of
at most five layers. Formula (1) represents each balanced term by four
linear lower bounds. The quadratic pair capacity is represented by an
exact finite table, and the capacity $c_i\leq n_i n_{i+1}$ by an integer
product. A graph necessarily supplies a feasible assignment. The converse
is not asserted: separately balanced degree sums need not be realized by
the same graph. A feasible assignment is therefore an obstruction to this
particular scalar exclusion argument, not a graph construction. The
program checks any returned assignment directly by integer arithmetic.
An infeasible solver report alone is not a formal exclusion certificate.

\paragraph{Attribution.}
The pair-counting principle and balanced minimum are the elementary
arguments already recorded in \texttt{layer\_pair\_bound.tex}; no novelty
is claimed for them. The additional step is to impose the corresponding
inequalities simultaneously on consecutive unions of distance layers.
The Moore interval and endpoint bounds are attributed in the retained
source proof to Cambie, Goedgebeur, Jooken and Van den Eede,
\emph{On the order-diameter ratio of girth-diameter cages},
arXiv:2511.21144v3, Propositions 1--2. This note changes no accepted
enumeration, source disposition, formal proof or compression measurement.
\end{document}
